% Zone „Das kennst du schon“ – aus bank/trigonometrische-funktionen/zone.jsonl (zusammenbau v0.1)
\einheitenkopf[zone][Kennst du schon]{Das kennst du schon}
\setcounter{aufgabe}{0}

%% TODO Titel: Ich-kann-Satz fehlt in der Bank (Platzhalter: Kettenname) – Nr. 1
%% TODO Anweisung: ein Satz über der ersten Teilaufgabe fehlt in der Bank – Nr. 1
\begin{aufgabe}{Sinus, Kosinus und Tangens als Seitenverhältnisse im rechtwinkligen Dreieck, Wert zu einem Winkel mit dem Taschenrechner im Grad-Modus}
\begin{teile}
\teil Gegenkathete 2 cm, Hypotenuse 5 cm – wie groß ist sin $\alpha$? sin $\alpha = $ \leerfeld
\teil sin 37° mit dem Taschenrechner im Grad-Modus – auf zwei Stellen? \leerfeld
\end{teile}
\end{aufgabe}

%% TODO Titel: Ich-kann-Satz fehlt in der Bank (Platzhalter: Kettenname) – Nr. 2
%% TODO Anweisung: ein Satz über der ersten Teilaufgabe fehlt in der Bank – Nr. 2
\begin{aufgabe}{Taschenrechner}
\begin{teile}
\teil Winkel in Grad – welchen Modus stellst du ein?\\ \kreuz{DEG}\\ \kreuz{RAD}
\teil sin $\alpha = 0{,}42$ – wie groß ist $\alpha$? Nutze die Umkehrtaste und runde auf eine Stelle. $\alpha \approx$ \leerfeld[°]
\end{teile}
\end{aufgabe}

%% TODO Titel: Ich-kann-Satz fehlt in der Bank (Platzhalter: Kettenname) – Nr. 3
%% TODO Anweisung: ein Satz über der ersten Teilaufgabe fehlt in der Bank – Nr. 3
\begin{aufgabe}{Koordinatensystem mit vier Quadranten, Punkte eintragen und ablesen, Achsen benennen und einteilen}
\begin{teile}
\teil Lies die Koordinaten von A ab.

\begin{ksys}[xmin=-4,xmax=4,ymin=-4,ymax=4,ablesen]\punkt{3}{2}{A}\end{ksys}
A(\leerfeld | \leerfeld)
\teil Lies die Koordinaten von C ab.

\begin{ksys}[xmin=-4,xmax=4,ymin=-4,ymax=4,xstep=0.5,ystep=0.5,ablesen]\punkt{-1.5}{-2.5}{C}\end{ksys}
C(\leerfeld | \leerfeld)
\end{teile}
\end{aufgabe}

%% TODO Titel: Ich-kann-Satz fehlt in der Bank (Platzhalter: Kettenname) – Nr. 4
%% TODO Anweisung: ein Satz über der ersten Teilaufgabe fehlt in der Bank – Nr. 4
\begin{aufgabe}{Kreis mit Mittelpunkt und Radius zeichnen, Umfang mit Pi berechnen, Mittelpunktswinkel und Kreisbogen als Anteil des Vollkreises}
\begin{teile}
\teil Radius 3 cm – wie lang ist der Umfang? Gib ihn mit $\pi$ und gerundet an. $u = $ \leerfeld[cm]
\teil Radius 4 cm, Mittelpunktswinkel 45° – wie lang ist der Kreisbogen? Runde auf zwei Stellen. $b \approx$ \leerfeld[cm]
\end{teile}
\end{aufgabe}

%% TODO Titel: Ich-kann-Satz fehlt in der Bank (Platzhalter: Kettenname) – Nr. 5
%% TODO Anweisung: ein Satz über der ersten Teilaufgabe fehlt in der Bank – Nr. 5
\begin{aufgabe}{Winkel messen, zeichnen und benennen, Vollwinkel, gestreckter und rechter Winkel, Winkel über neunzig Grad als stumpfe Winkel}
\begin{teile}
\teil Wie groß ist ein gestreckter Winkel? \leerfeld[°]
\teil Zeichne am Strahl einen Winkel von 135° ein.

\winkelstrahl
\end{teile}
\end{aufgabe}

%% TODO Titel: Ich-kann-Satz fehlt in der Bank (Platzhalter: Kettenname) – Nr. 6
%% TODO Anweisung: ein Satz über der ersten Teilaufgabe fehlt in der Bank – Nr. 6
\begin{aufgabe}{Bruchteile eines Vollkreises als Bruch und als Dezimalzahl (halb, viertel, drittel)}
\begin{teile}
\teil Ein Viertel des Vollkreises – wie viel Grad? \leerfeld[°]
\teil 72° – welcher Anteil des Vollkreises? Gib ihn als Bruch und als Dezimalzahl an. \leerfeld
\end{teile}
\end{aufgabe}

%% TODO Titel: Ich-kann-Satz fehlt in der Bank (Platzhalter: Kettenname) – Nr. 7
%% TODO Anweisung: ein Satz über der ersten Teilaufgabe fehlt in der Bank – Nr. 7
\begin{aufgabe}{Wertetabelle anlegen, Wertepaare als Punkte eintragen, den Graphen als durchgehende Linie zeichnen, Funktionswert und Argument ablesen, Punktprobe}
\begin{teile}
\teil $y = 2x$ – fülle die Wertetabelle aus.

\wertetabelle{x}{y}{0,1,2,3}
\teil $y = x^2 - 1$ – fülle die Wertetabelle aus.

\wertetabelle{x}{y}{-2,-1,0,1,2}
\end{teile}
\end{aufgabe}

%% TODO Titel: Ich-kann-Satz fehlt in der Bank (Platzhalter: Kettenname) – Nr. 8
%% TODO Anweisung: ein Satz über der ersten Teilaufgabe fehlt in der Bank – Nr. 8
\begin{aufgabe}{Achseneinteilung wählen, wenn beide Achsen verschiedene Einheiten und Schrittweiten haben}
\begin{teile}
\teil Rechtsachse: 0 bis 60 Minuten auf 6 cm – wie viele Minuten je Zentimeter? \leerfeld[min]
\teil Rechtsachse: 0 bis 2,5 Stunden auf 10 cm – wie viel je Zentimeter? \leerfeld[h]
\end{teile}
\end{aufgabe}

%% TODO Titel: Ich-kann-Satz fehlt in der Bank (Platzhalter: Kettenname) – Nr. 9
%% TODO Anweisung: ein Satz über der ersten Teilaufgabe fehlt in der Bank – Nr. 9
\begin{aufgabe}{Achsensymmetrisch und punktsymmetrisch als Wörter, Symmetrieachse und Symmetriezentrum am Graphen}
\begin{teile}
\teil $y = x^2$ – welche Symmetrie?\\ \kreuz{achsensymmetrisch zur Hochachse}\\ \kreuz{punktsymmetrisch zum Ursprung}
\teil $y = (x - 2)^2$ – wo liegt die Symmetrieachse? $x = $ \leerfeld
\end{teile}
\end{aufgabe}

%% TODO Titel: Ich-kann-Satz fehlt in der Bank (Platzhalter: Kettenname) – Nr. 10
%% TODO Anweisung: ein Satz über der ersten Teilaufgabe fehlt in der Bank – Nr. 10
\begin{aufgabe}{Parameter in einer Funktionsgleichung erkennen und ihre Wirkung auf den Graphen beschreiben (Streckung, Stauchung, Verschiebung), Graph zu Gleichung zuordnen}
\begin{teile}
\teil Welche Parabel ist schmaler?\\ \kreuz{$y = 3x^2$}\\ \kreuz{$y = x^2$}
\teil $y = -0{,}25x^2$ – wie wirkt der Faktor vor $x^2$? \leerfeld
\end{teile}
\end{aufgabe}

%% TODO Titel: Ich-kann-Satz fehlt in der Bank (Platzhalter: Kettenname) – Nr. 11
%% TODO Anweisung: ein Satz über der ersten Teilaufgabe fehlt in der Bank – Nr. 11
\begin{aufgabe}{Größtwert, Kleinstwert und Spannweite einer Tabelle entnehmen, Mittelwert bilden}
\begin{teile}
\teil Temperaturen 4, 9, 7, 11, 6 (in °C) – Größtwert und Kleinstwert? \leerfeld; \leerfeld
\teil −2, 5, 3, −1, 4 (in °C) – wie groß ist der Mittelwert? \leerfeld[°C]
\end{teile}
\end{aufgabe}

%% TODO Titel: Ich-kann-Satz fehlt in der Bank (Platzhalter: Kettenname) – Nr. 12
%% TODO Anweisung: ein Satz über der ersten Teilaufgabe fehlt in der Bank – Nr. 12
\begin{aufgabe}{Taschenrechner – Fehler finden}
\begin{teile}
\teil Tom rechnet 8 · sin 35° und erhält −3,43. Finde den Fehler und rechne richtig.
\end{teile}
\end{aufgabe}

%% TODO Titel: Ich-kann-Satz fehlt in der Bank (Platzhalter: Kettenname) – Nr. 13
%% TODO Anweisung: ein Satz über der ersten Teilaufgabe fehlt in der Bank – Nr. 13
\begin{aufgabe}{Taschenrechner – selbst rechnen}
\begin{teile}
\teil 6 · cos 48° – auf zwei Stellen? Prüfe vorher den Modus. \leerfeld
\end{teile}
\end{aufgabe}
